The approach to obtaining the classification is to take the 1D
revival set $ k,m \in \{2, 3, 4, 5, 6, 8, 10\}$ established in Theorem~1, together with the previously found 1D parameter values for all known revivals in 
Dukes~\cite{Dukes2014} as a starting point.

We then assemble the 1D
walk operator $U_k = S_k(I_k \otimes C_k)$ 
using the coin of Eq.~(\ref{eq:1d_coin}) and verify each tabulated
revival by direct matrix exponentiation, checking that $U_k^N$ equals
the identity within a tolerance value of $10^{-8}$. 

The 2D classification then follows from the tensor-product
factorisation $U_{k \times m} = U_k \otimes U_m$ of
Eq.~(\ref{eq:factorised_unitary}). For each pair $(k, m)$ with
$k, m \in \{2, 3, 4, 5, 6, 8, 10\}$, we enumerate the periods $N$ shared
by the two 1D revival sets $\mathcal{N}_k$ and $\mathcal{N}_m$ and, for
each such $N$, record every combination of 1D parameters
$(\rho_x, \delta_x) \in \mathcal{R}_k(N)$ and
$(\rho_y, \delta_y) \in \mathcal{R}_m(N)$ that produces the revival,

Because the 1D input is complete and the 2D pairing is
exhaustive over the finite revival set, the resulting classification revivals
with different values for $\rho_x, \rho_y, \delta_x, \delts_y$ is obtained. The full classification is available at \cite{full_revivals_references}.
